\section{Theoretical Analysis}
\label{app:dfb-theory}
\label{app:theory}
\label{app:bottleneck-proof}

\subsection{Proof of Lemma~\ref{lem:linear-action-difference}}

\paragraph{Setup.}
Fix $z$ and write $\pi=\pi_z$.
Let $\rho$ be the reference probability distribution, $r$ a bounded measurable reward, and $B\in L^2(\rho;\mathbb R^d)$ fixed. 
At a bottleneck entry state $s$, choose $a_+\in\mathcal A_+(s)$ and $a_-\in\mathcal A_-(s)$ such that $Q_r^\pi(s,a_+)>Q_r^\pi(s,a_-)$.
For each $a\in\{a_+,a_-\}$, write
\[
    M_a:=M^\pi(s,a,\cdot),
    \qquad
    M_a(ds^+)=m_a(s^+)\,\rho(ds^+),
    \qquad
    m_a\in L^2(\rho).
\]
The forward vectors minimize their squared errors independently:
\[
    F^*(s,a,z)
    \in\operatorname*{arg\,min}_{f\in\mathbb R^d}
    \int_{\mathcal S}
    \bigl[m_a(s^+)-f^\top B(s^+)\bigr]^2\,\rho(ds^+).
\]
Define the fitted densities and values by
\[
    \widehat m_a(s^+):=F^*(s,a,z)^\top B(s^+),
    \qquad
    \widehat Q_r^\pi(s,a)
    :=\frac{1}{1-\gamma}
    \int_{\mathcal S}r(s^+)\widehat m_a(s^+)\,\rho(ds^+).
\]

\begin{proof}
\noindent\textbf{Step 1: Characterize the optimal fit.}
With $B$ fixed, fitting retains the component of each successor density that lies in the span of the backward features.
To see this, differentiating the squared-error objective gives
\[
    G_B F^*(s,a,z)
    =\int_{\mathcal S}B(s^+)\,M_a(ds^+),
    \qquad
    G_B:=\int_{\mathcal S}
    B(s^+)B(s^+)^\top\,\rho(ds^+).
\]
Equivalently, for every $j=1,\ldots,d$,
\[
    \int_{\mathcal S}
    B_j(s^+)\bigl[m_a(s^+)-\widehat m_a(s^+)\bigr]
    \,\rho(ds^+)=0.
\]
Thus the residual is orthogonal to the backward feature span.
Writing $\Pi_B$ for the $L^2(\rho)$ orthogonal projection onto $\operatorname{span}\{B_1,\ldots,B_d\}$, we obtain
\[
    \widehat m_a=\Pi_Bm_a
    \qquad \rho\text{-almost everywhere}.
\]
The fitted density is therefore unique, even if $G_B$ is singular and the minimizing forward vector is not unique.

\noindent\textbf{Step 2: Identify when the action distinction is lost.}
The two fits coincide exactly when the backward features have the same expectations under the two SMs.
Let
\[
    v:=F^*(s,a_+,z)-F^*(s,a_-,z).
\]
Subtracting the optimality equations yields
\[
    G_Bv
    =\int_{\mathcal S}B(s^+)\,M_{a_+}(ds^+)
     -\int_{\mathcal S}B(s^+)\,M_{a_-}(ds^+).
\]
If these feature expectations agree, then $G_Bv=0$, and hence
\[
    \int_{\mathcal S}
    \bigl[\widehat m_{a_+}(s^+)-\widehat m_{a_-}(s^+)\bigr]^2
    \,\rho(ds^+)
    =v^\top G_Bv=0.
\]
The fitted densities therefore coincide $\rho$-almost everywhere.
Conversely, if they coincide, then
\[
    G_Bv
    =\int_{\mathcal S}B(s^+)
    \bigl[\widehat m_{a_+}(s^+)-\widehat m_{a_-}(s^+)\bigr]
    \,\rho(ds^+)=0,
\]
so their feature expectations agree.
This proves
\[
    \mathbb E_{M_{a_+}}[B]=\mathbb E_{M_{a_-}}[B]
    \quad\Longleftrightarrow\quad
    \widehat m_{a_+}=\widehat m_{a_-}
    \quad \rho\text{-almost everywhere}.
\]

\noindent\textbf{Step 3: Relate the lost distinction to action values.}
The same projection can be applied to the reward: evaluating the fitted density against $r$ is equivalent to evaluating the exact density against $\Pi_Br$.
Indeed, orthogonality gives
\[
\begin{aligned}
    (1-\gamma)\widehat Q_r^\pi(s,a)
    &=\int_{\mathcal S}
      r(s^+)\widehat m_a(s^+)\,\rho(ds^+)\\
    &=\int_{\mathcal S}
      (\Pi_Br)(s^+)\widehat m_a(s^+)\,\rho(ds^+)\\
    &=\int_{\mathcal S}
      (\Pi_Br)(s^+)m_a(s^+)\,\rho(ds^+).
\end{aligned}
\]
The second equality uses that $\widehat m_a$ lies in the backward feature span; the third uses that $\Pi_Br$ does.

Under the equivalent conditions of Step 2, the fitted densities coincide, so
\[
    \widehat Q_r^\pi(s,a_+)-\widehat Q_r^\pi(s,a_-)
    =\frac{1}{1-\gamma}
    \int_{\mathcal S}r(s^+)
    \bigl[\widehat m_{a_+}(s^+)-\widehat m_{a_-}(s^+)\bigr]
    \,\rho(ds^+)=0.
\]
For the chosen action pair, however, \eqref{eq:sm-value} gives
\[
    Q_r^\pi(s,a_+)-Q_r^\pi(s,a_-)
    =\frac{1}{1-\gamma}
    \int_{\mathcal S}r(s^+)
    \bigl[m_{a_+}(s^+)-m_{a_-}(s^+)\bigr]
    \,\rho(ds^+)>0.
\]
Thus fitting discards the component of the successor densities that carries their true action-value distinction.
\end{proof}

\subsection{Proof of Proposition~\ref{prop:bottleneck-preservation}}

\paragraph{Setup.}
Both the normalized SMs $M^\pi(s,a,\cdot)$ and the Density-FB predictions $\widehat M_\theta(s,a,z,\cdot)$ are probability measures.
Define the fitted values by
\[
    \widehat Q_r^\pi(s,a) :=\frac{1}{1-\gamma} \int_{\mathcal S} r(s^+)\,\widehat M_\theta(s,a,z,ds^+).
\]
The assumed positive gap implies that $r$ is non-constant, so $0<R<\infty$.

\begin{proof}
\noindent\textbf{Step 1: Convert distribution error into value error.}
For bounded rewards, closeness of the successor distributions controls the error in their predicted returns.
The bounded-reward expectation bound and Pinsker's inequality give, for each $a\in\{a_+,a_-\}$,
\[
    \left| \widehat Q_r^\pi(s,a)-Q_r^\pi(s,a) \right| \le \frac{R}{1-\gamma} \sqrt{ \frac12 D_{\mathrm{KL}}\!\left( M^\pi(s,a,\cdot) \,\middle\|\, \widehat M_\theta(s,a,z,\cdot) \right) }.
\]
Here KL divergence uses natural-log.

\noindent\textbf{Step 2: Compare the errors with the true gap.}
The stated KL threshold makes each value error smaller than half the true action-value gap:
\[
    \left| \widehat Q_r^\pi(s,a)-Q_r^\pi(s,a) \right| < \frac{R}{1-\gamma} \sqrt{\frac{(1-\gamma)^2\delta^2}{4R^2}} =\frac{\delta}{2}, \qquad a\in\{a_+,a_-\}.
\]

\noindent\textbf{Step 3: Preserve the action ordering.}
Even if both errors reduce the predicted gap, their combined magnitude is insufficient to erase it:
\[
\begin{aligned}
    &\widehat Q_r^\pi(s,a_+)-\widehat Q_r^\pi(s,a_-)\\
    &\quad\ge
    \delta -\left|\widehat Q_r^\pi(s,a_+)-Q_r^\pi(s,a_+)\right| -\left|\widehat Q_r^\pi(s,a_-)-Q_r^\pi(s,a_-)\right|\\
    &\quad> \delta-\frac{\delta}{2}-\frac{\delta}{2} =0.
\end{aligned}
\]
\end{proof}


\subsection{Convergence under exact fitting}
\label{app:exact-fitting-convergence}

We adapt the exact-expectation convergence argument of TD InfoNCE~\citep[Appendix~A]{zheng-iclr24a} and establish a KL contraction that connects the fitted updates to \cref{prop:bottleneck-preservation}.

\paragraph{Setup.}
Assume finite state and action spaces, $0<\gamma<1$, a reference distribution $\rho$ with positive mass on every state, and a fixed continuation policy $\pi=\pi_z$.
All expectations and normalizing constants are evaluated exactly.
Write
\[
    \widehat M_k(s,a,\cdot) :=\widehat M_{\theta_k}(s,a,z,\cdot),
\]
and initialize $\widehat M_0(s,a,\cdot)$ with positive mass on every state for every $(s,a)$.

At iteration $k$, form the Bellman targets
\[
    T_k(s,a,\cdot) :=(1-\gamma)P(\cdot\mid s,a) +\gamma \mathbb E_{\substack{ s'\sim P(\cdot\mid s,a)\\ a'\sim\pi(\cdot\mid s') }} \left[\widehat M_k(s',a',\cdot)\right].
\]
Assume that these targets are jointly representable by \name{} at the chosen feature dimension.
With the targets held fixed, globally minimize the weighted sum of categorical fitting losses, assigning positive weight to every state-action pair.
Both the forward and backward parameters may change during fitting.
Then
\[
    \max_{s,a} D_{\mathrm{KL}}\!\left( M^\pi(s,a,\cdot) \,\middle\|\, \widehat M_k(s,a,\cdot) \right) \longrightarrow 0.
\]
Consequently, for any fixed bounded reward and any bottleneck action pair with a strictly positive true value gap, the fitted ordering is correct at every sufficiently large iteration.

\begin{proof}
\noindent\textbf{Step 1: Recover the Bellman target.}
Exact cross-entropy fitting recovers a representable target distribution.
Writing
\[
    \widehat M_\theta(s,a,z,ds^+)
    =h_\theta(s,a,z,s^+)\,\rho(ds^+),
\]
the loss for each state-action pair satisfies
\[
\begin{aligned}
    &-\int_{\mathcal S}
      \log h_\theta(s,a,z,s^+)\,T_k(s,a,ds^+)\\
    &\qquad=
    D_{\mathrm{KL}}\!\left(
        T_k(s,a,\cdot)
        \,\middle\|\,
        \widehat M_\theta(s,a,z,\cdot)
    \right)
    +\text{constant},
\end{aligned}
\]
where the constant is independent of $\theta$.
Joint realizability allows all KL terms to vanish simultaneously.
Since every fitting weight is positive, every global minimizer therefore satisfies
\[
    \widehat M_{k+1}(s,a,\cdot)=T_k(s,a,\cdot)
    \qquad\text{for every }(s,a).
\]

\noindent\textbf{Step 2: Contract the KL error.}
The true SM and the fitted update have the same immediate-transition component.
Their discrepancy comes only from the continuation predictions.
Define
\[
    K_k:=\max_{s,a}
    D_{\mathrm{KL}}\!\left(
        M^\pi(s,a,\cdot)
        \,\middle\|\,
        \widehat M_k(s,a,\cdot)
    \right).
\]
Using \eqref{eq:sm-bellman}, Step 1, and joint convexity of KL,
\[
\begin{aligned}
    &D_{\mathrm{KL}}\!\left(
        M^\pi(s,a,\cdot)
        \,\middle\|\,
        \widehat M_{k+1}(s,a,\cdot)
    \right)\\
    &\quad\le
    \gamma
    \mathbb E_{\substack{
        s'\sim P(\cdot\mid s,a)\\
        a'\sim\pi(\cdot\mid s')
    }}
    \left[
        D_{\mathrm{KL}}\!\left(
            M^\pi(s',a',\cdot)
            \,\middle\|\,
            \widehat M_k(s',a',\cdot)
        \right)
    \right]\\
    &\quad\le\gamma K_k.
\end{aligned}
\]
Taking the maximum over $(s,a)$ gives
\[
    K_{k+1}\le\gamma K_k.
\]

\noindent\textbf{Step 3: Reach the preservation threshold.}
Positive initialization and finite state and action spaces ensure $K_0<\infty$.
Hence
\[
    K_k\le\gamma^kK_0\longrightarrow0.
\]
Fix a bounded reward $r$ and a bottleneck action pair $a_+\in\mathcal A_+(s)$, $a_-\in\mathcal A_-(s)$ with
\[
    \delta:=Q_r^\pi(s,a_+)-Q_r^\pi(s,a_-)>0.
\]
Let
\[
    R:=\sup_{\mathcal S}r-\inf_{\mathcal S}r,
    \qquad
    \varepsilon:=\frac{(1-\gamma)^2\delta^2}{2R^2}.
\]
There exists $k_0$ such that, for every $k\ge k_0$,
\[
    D_{\mathrm{KL}}\!\left(
        M^\pi(s,a,\cdot)
        \,\middle\|\,
        \widehat M_k(s,a,\cdot)
    \right)
    \le K_k<\varepsilon,
    \qquad a\in\{a_+,a_-\}.
\]
Thus \cref{prop:bottleneck-preservation} gives
\[
    \widehat Q_{k,r}^\pi(s,a_+)>
    \widehat Q_{k,r}^\pi(s,a_-),
    \qquad k\ge k_0,
\]
where $\widehat Q_{k,r}^\pi$ integrates $r$ against $\widehat M_k$ as in \eqref{eq:sm-value}.
\end{proof}


\subsection{Successor-feature interpretation of the training reward}
\label{app:reward-successor-features}

\paragraph{Setup.}
Fix $z$, $\pi=\pi_z$, and a bounded measurable feature map $B_\theta:\mathcal S\to\mathbb R^d$.
Consider the training reward
\[
    r_z(s^+):=\frac{B_\theta(s^+)^\top z}{\sqrt d}.
\]
With these reward features held fixed, the value of $r_z$ has the standard successor-feature representation~\citep{barreto-neurips17}.

\begin{proof}
\noindent\textbf{Step 1: Express successor features using the SM.}
Successor features record the expected discounted sum of future features.
By the definition of the normalized SM,
\[
\begin{aligned}
    \psi^\pi(s,a)
    &:=
    \mathbb E_\pi\!\left[
        \sum_{t\ge0}\gamma^t
        \frac{B_\theta(s_{t+1})}{\sqrt d}
        \,\middle|\,
        s_0=s,a_0=a
    \right]\\
    &=
    \frac{1}{1-\gamma}
    \int_{\mathcal S}
    \frac{B_\theta(s^+)}{\sqrt d}\,
    M^\pi(s,a,ds^+).
\end{aligned}
\]
Boundedness of $B_\theta$ justifies exchanging the expectation and the discounted sum.

\noindent\textbf{Step 2: Recover the reward value.}
Since the reward is linear in these features, its value is linear in their discounted occupancies:
\[
\begin{aligned}
    Q_{r_z}^\pi(s,a)
    &=
    \frac{1}{1-\gamma}
    \int_{\mathcal S}
    \frac{B_\theta(s^+)^\top z}{\sqrt d}\,
    M^\pi(s,a,ds^+)\\
    &=\psi^\pi(s,a)^\top z.
\end{aligned}
\]

\noindent\textbf{Step 3: Identify the features estimated by the readout.}
For candidate future states $s_j^+$, the probabilities $p_j(s,a,z)$ in \eqref{eq:density-candidates} define the estimate
\[
    \widehat\psi_\theta(s,a,z)
    :=
    \frac{1}{1-\gamma}
    \sum_j p_j(s,a,z)
    \frac{B_\theta(s_j^+)}{\sqrt d}.
\]
Substituting $r_z$ into \eqref{eq:density-value} gives
\[
\begin{aligned}
    \widehat Q_{\theta,r_z}(s,a,z)
    &=
    \frac{1}{1-\gamma}
    \sum_j p_j(s,a,z)\,r_z(s_j^+)\\
    &=\widehat\psi_\theta(s,a,z)^\top z.
\end{aligned}
\]
Thus the readout estimates successor features by averaging future-state features under the predicted probabilities, then applies the reward weights $z$.
\end{proof}

For the implemented centered reward, the same identities hold with $B_\theta$ replaced by $B_\xi-\hat\mu$, holding $B_\xi$ and $\hat\mu$ fixed.